\documentclass[11pt]{article}
\usepackage[margin=1in]{geometry}
\usepackage[T1]{fontenc}
\usepackage[utf8]{inputenc}
\usepackage{amsmath,amssymb,amsthm}
\usepackage{enumitem}

\title{ICPC World Finals 2020\\N. What’s Our Vector, Victor?}
\author{}
\date{}

\begin{document}
\maketitle

\section*{Problem Summary}

We are given $n$ known vectors $p_1,\dots,p_n \in \mathbb{R}^d$ and the Euclidean distances from an
unknown vector $x$ to each of them. We must output any vector $x$ consistent with all those distances.

The hidden test data is generated from an actual vector, so a solution always exists.

\section*{Key Observations}

\begin{itemize}[leftmargin=*]
    \item Let $y = x - p_1$. Then $\|y\| = e_1$.
    \item For every $i \ge 2$, let $v_i = p_i - p_1$. Since $\|x-p_i\| = e_i$, we get
    \[
        \|y - v_i\|^2 = e_i^2.
    \]
    Subtracting $\|y\|^2 = e_1^2$ removes the quadratic term:
    \[
        2\,v_i \cdot y = \|v_i\|^2 + e_1^2 - e_i^2.
    \]
    So all constraints except the first become linear equations in $y$.
    \item The linear equations define an affine subspace. We need any point in that subspace whose norm is
    exactly $e_1$.
    \item If we compute a basis of the nullspace of the linear system, then the point of minimum norm in the
    affine subspace is obtained by projecting any particular solution away from the nullspace.
    \item The difference between $e_1^2$ and that minimum norm tells us how much freedom remains in a
    nullspace direction.
\end{itemize}

\section*{Algorithm}

\begin{enumerate}[leftmargin=*]
    \item Build the linear system
    \[
        v_i \cdot y = \frac{\|v_i\|^2 + e_1^2 - e_i^2}{2}
        \qquad (i = 2,\dots,n).
    \]
    \item Run Gaussian elimination on the augmented matrix to obtain:
    \begin{itemize}[leftmargin=*]
        \item one particular solution $y_0$ (set all free variables to $0$),
        \item one basis vector for each free variable in the nullspace.
    \end{itemize}
    \item Orthonormalize the nullspace basis with Gram--Schmidt.
    \item Project $y_0$ away from the nullspace to get the minimum-norm solution $y_{\min}$.
    \item If $\|y_{\min}\| = e_1$, we are done.
    Otherwise, add
    \[
        \sqrt{e_1^2 - \|y_{\min}\|^2}
    \]
    times any unit nullspace vector. This stays inside the affine subspace and fixes the norm.
    \item Output $x = p_1 + y$.
\end{enumerate}

\section*{Correctness Proof}

We prove that the algorithm returns the correct answer.

\paragraph{Lemma 1.}
For every $i \ge 2$, the constraint $\|x-p_i\| = e_i$ is equivalent to the linear equation
\[
    v_i \cdot y = \frac{\|v_i\|^2 + e_1^2 - e_i^2}{2},
\]
where $y=x-p_1$ and $v_i=p_i-p_1$.

\paragraph{Proof.}
We have $\|y\|=\|x-p_1\|=e_1$ and $\|y-v_i\|=\|x-p_i\|=e_i$.
Expanding $\|y-v_i\|^2=e_i^2$ gives
\[
    \|y\|^2 - 2v_i \cdot y + \|v_i\|^2 = e_i^2.
\]
Substituting $\|y\|^2=e_1^2$ and rearranging yields the stated linear equation. \qed

\paragraph{Lemma 2.}
Every solution of the original problem is exactly a vector of the form $x=p_1+y$, where $y$ belongs to
the affine subspace defined by the linear equations from Lemma 1 and also satisfies $\|y\|=e_1$.

\paragraph{Proof.}
By Lemma 1, satisfying all distances except the first is equivalent to belonging to the affine subspace.
The first distance is exactly $\|x-p_1\|=\|y\|=e_1$. Therefore the original system and the stated affine
subspace plus norm condition are equivalent. \qed

\paragraph{Lemma 3.}
Let $A$ be the linear system from Lemma 1. After Gaussian elimination, the algorithm constructs:
\begin{itemize}[leftmargin=*]
    \item a particular solution $y_0$ of $Ay=b$,
    \item a basis of the nullspace $\mathcal{N}=\{z:Az=0\}$.
\end{itemize}
Then every solution of $Ay=b$ is of the form $y_0+z$ with $z \in \mathcal{N}$.

\paragraph{Proof.}
This is the standard characterization of solutions to a consistent linear system after row reduction. Free
variables generate the nullspace, and setting them all to zero gives one particular solution. \qed

\paragraph{Lemma 4.}
The vector $y_{\min}$ produced by projecting $y_0$ away from the nullspace is the minimum-norm solution
of $Ay=b$.

\paragraph{Proof.}
After orthonormalizing the nullspace basis, subtracting from $y_0$ its orthogonal projection onto the
nullspace gives a vector $y_{\min}$ that is orthogonal to the nullspace.
Any other solution has the form $y_{\min}+z$ with $z \in \mathcal{N}$ by Lemma 3, and orthogonality gives
\[
    \|y_{\min}+z\|^2 = \|y_{\min}\|^2 + \|z\|^2 \ge \|y_{\min}\|^2.
\]
So $y_{\min}$ is exactly the minimum-norm solution. \qed

\paragraph{Lemma 5.}
If the affine subspace contains any vector of norm $e_1$, then the algorithm constructs one.

\paragraph{Proof.}
By Lemma 4, every solution has norm at least $\|y_{\min}\|$.
If $\|y_{\min}\|=e_1$, the algorithm already has a valid vector.
Otherwise, because a valid vector exists, the affine subspace must have a nontrivial nullspace direction.
Moving along any unit nullspace vector keeps us inside the affine subspace, and by orthogonality to
$y_{\min}$ the norm becomes
\[
    \|y_{\min} + t u\|^2 = \|y_{\min}\|^2 + t^2.
\]
Choosing $t = \sqrt{e_1^2 - \|y_{\min}\|^2}$ gives norm exactly $e_1$. \qed

\paragraph{Theorem.}
The algorithm outputs the correct answer for every valid input.

\paragraph{Proof.}
By Lemma 2, it is enough to find a vector $y$ in the affine subspace with norm $e_1$.
Lemmas 3, 4, and 5 show that the algorithm constructs such a vector, and therefore $x=p_1+y$ satisfies
all required distances. \qed

\section*{Complexity Analysis}

Gaussian elimination on at most $499$ equations and $500$ variables takes $O(nd^2)$ time here, and the
nullspace orthonormalization costs $O(d^3)$ in the worst case. With $d,n \le 500$, this easily fits.
The memory usage is $O(d^2)$.

\section*{Implementation Notes}

\begin{itemize}[leftmargin=*]
    \item Using `long double` keeps the numerical error comfortably below the $10^{-5}$ tolerance.
    \item The generated test data is consistent, so we do not need to handle infeasible systems.
    \item Small negative values caused by floating-point roundoff are clamped to zero before taking square
    roots.
\end{itemize}

\end{document}
